\FloatBarrier\section{Cuaderno integrado local y Google Colab}

\subsection{Identidad del estudio y código compartido}
El cuaderno \texttt{HAM10000\_Grupo09.ipynb} adopta el mismo título que este informe: \textit{Evaluación Comparativa de Tres Arquitecturas CNN para Clasificación de Lesiones Cutáneas Pigmentadas mediante Dermatoscopia (HAM10000)}. HAM10000 identifica el conjunto de imágenes, no una ubicación del estudio ni una herramienta clínica. Se conserva el alcance académico, sin despliegue.

El notebook integra el flujo completo mediante el paquete común \texttt{python/ham10000}, la CLI \texttt{ejecutar.py}, \texttt{configs/base.json} y el manifiesto fijo. No mantiene otra implementación del entrenamiento. Sus 22 celdas cubren controles, inicialización, dependencias, inventario del código, integridad de datos, protocolo, corridas existentes, entrenamiento/reanudación, evaluación, métricas, figuras, discusión, exportación y diagnóstico. \texttt{generar\_notebook.py} permite mantener su fuente revisable; las utilidades de \texttt{runtime.py} se integran en el propio cuaderno.

\subsection{Detección de entorno y resolución del error local}
La importación incondicional de \texttt{google.colab} produce \texttt{ModuleNotFoundError} al ejecutar un cuaderno de nube desde VS Code o Jupyter ordinario. Instalar ese paquete no habilita el montaje de Drive fuera del backend compatible. La versión corregida detecta el runtime alojado mediante sus variables de entorno y solo importa/monta Drive dentro de esa rama \cite{colabfaq}. Localmente descubre la raíz del proyecto y los originales, o permite indicar rutas explícitas, e imprime \texttt{sys.executable}, rutas y versiones reales. La inspección no exige CUDA; el entrenamiento base sí queda bloqueado si no hay GPU disponible.

\subsection{Paquete portable, extracción e integridad}
El script \texttt{empaquetar.py} genera \texttt{proyecto\_ham10000\_colab.zip}; informa su tamaño y comprueba CRC. Incluye el paquete Python completo, requisitos CUDA/comunes/Jupyter, configuración, manifiestos, fuentes LaTeX, tablas, figuras, PDF y presentación. También incluye \texttt{resultados/comparacion.csv} como \textbf{referencia del experimento local publicado}; no es evidencia de una ejecución nueva. Excluye imágenes, carpetas exportadas, entorno virtual, checkpoints, \texttt{runs/} y resultados pesados.

En Colab el usuario sube manualmente ese ZIP y \texttt{HAM10000.zip} a \texttt{Mi unidad/Grupo09HAM/}. El cuaderno autoriza Drive, copia y extrae en \texttt{/content/ham10000\_sesion/} para evitar miles de lecturas remotas. Comprueba SHA-256 de las copias y CRC al leer, rechaza rutas inseguras, enlaces, colisiones incompatibles y espacio insuficiente; publica el directorio solo al terminar. Una extracción no acreditada no se sobrescribe. Después verifica el digest canónico de \texttt{split.csv}, los conteos, clases y fugas por lesión/grupo/hash; por defecto también rehashea los 10\,015 originales, sin volver a sortear particiones.

\subsection{Persistencia y controles seguros}
\texttt{ENTRENAR}, \texttt{EVALUAR}, \texttt{CONFIRMAR\_TEST}, instalación, regeneración de figuras y exportación comienzan desactivados. La configuración de entrenamiento se conserva íntegra, incluido \texttt{num\_workers=4}; ya no se genera automáticamente una variante con dos workers. Reanudar exige arquitectura, configuración completa y split idénticos. Se continúa desde la última época completa de \texttt{last.pt}; una corrida terminada no se reescribe. Para otro protocolo se requiere otro JSON y raíz de runs explícitos, sin mezclarlo con el experimento publicado.

En Colab los checkpoints, historiales y métricas se escriben en Drive; los originales se leen desde el almacenamiento temporal. Al reiniciar se reconstruye \texttt{/content} y se utiliza el mismo protocolo. No se garantizan GPU gratuita, duración de sesión, identidad bit a bit entre dispositivos ni persistencia ante fallos del montaje o almacenamiento.

\begin{table}[htbp]\centering\small
\caption{Equivalencia del flujo y diferencias operativas entre ejecución local y Colab.}
\begin{tabularx}{\textwidth}{@{}>{\bfseries}p{3.5cm}>{\raggedright\arraybackslash}X>{\raggedright\arraybackslash}X@{}}
\toprule
\textbf{Aspecto} & \textbf{Local (VS Code/Jupyter)} & \textbf{Google Colab alojado} \\ \midrule
Código y protocolo & Paquete común; JSON original & Mismo paquete y JSON desde ZIP \\
Drive & No se importa ni monta & Montaje autorizado solo en backend compatible \\
Originales & Disco local, rutas explícitas o descubiertas & Extracción temporal en \texttt{/content} \\
Partición y workers & Split \texttt{583779fb...}; workers 4 & Mismo split y workers 4 \\
Persistencia & \texttt{runs/} local & \texttt{Grupo09HAM/runs} en Drive \\
Resultados nuevos & \texttt{resultados\_notebook/} & Misma carpeta bajo Drive \\
Evidencia acreditada & Entrenamientos reales y cuaderno en inspección & Rama simulada; no ejecución real acreditada \\
\bottomrule
\end{tabularx}
\end{table}

\subsection{Evaluación, resultados e integración documental}
Con evaluación activada se ejecuta primero \texttt{val\_clean}; solo tras confirmar explícitamente el protocolo se realizan las 21 evaluaciones finales de test (cinco condiciones y dos mitigaciones por CNN). Antes de reutilizar una carpeta, \texttt{artifacts.py} comprueba ambos archivos, procedencia, todas las muestras, etiquetas y probabilidades. Una carpeta incompleta detiene el flujo; no se omite silenciosamente. Las salidas nuevas se publican desde un directorio temporal al completar CSV y JSON. \texttt{-{}-device cpu} y \texttt{-{}-num-workers 0} permiten evaluar un checkpoint CUDA sin editarlo; las mediciones CPU se conservan separadas.

El cuaderno distingue métricas verificadas en los runs de la tabla de referencia local y no fabrica resultados. Expone Macro Precision/F1, exactitud balanceada, Recall de melanoma, AUC, reportes por clase, matrices, robustez, mitigación y eficiencia. Las figuras opcionales se generan a color y 300 DPI en \texttt{resultados\_notebook/}, sin sustituir cifras ni gráficos publicados. La consolidación permite seleccionar subcarpeta y rechaza duplicados o entornos de inferencia mezclados; tiempo y VRAM de entrenamiento requieren revisar además las sesiones registradas.

La exportación opcional conserva CSV/JSON, configuración, split y figuras con un manifiesto de procedencia y SHA-256; no incluye imágenes ni checkpoints. Actualizar informe, artículo y exposición con un estudio nuevo exige revisar deliberadamente evidencias y textos, compilar y solo después empaquetar. Se verificaron esquema nbformat, sintaxis, regresiones y ejecución local del cuaderno en modo de inspección, con artefactos experimentales sin cambios. La rama Colab se simuló sin autenticación externa: \textbf{no se acredita ejecución en una cuenta real de Google Colab}. Todas las métricas publicadas siguen correspondiendo a la GTX~1650 local.
